\section{Related Work}
\label{sec:related_work}

\subsection{Deep reinforcement learning for HVAC control}
\label{ssec:rw_drl}

Deep reinforcement learning has been applied to HVAC supervisory
control for the better part of a decade, starting with discrete-action
$Q$-learning controllers \citep{Wei2017DeepRL,Mnih2015Atari} and
moving rapidly to continuous-action policy-gradient methods such as
PPO and SAC \citep{Schulman2017PPO,Haarnoja2018SAC}. Recent reviews
catalogue dozens of such studies on platforms ranging from in-house
EnergyPlus wrappers through Sinergym \citep{Sinergym2021} to the
BOPTEST suite \citep{Blum2021BOPTEST}; surveys consistently report
that the bottleneck to deployment is not algorithmic but environmental
and statistical \citep{Mason2019RLReview,AlSayed2024Deployment}. Two
recurring weaknesses are particularly relevant here. First, the vast
majority of DRL-HVAC studies report a single seed or at most three,
making energy-comfort comparisons fragile under environment
stochasticity \citep{Henderson2018Reproducibility,Agarwal2021RLPrecipice,
Hedayat2025FiftySeed}. Second, the choice of \emph{which} surrogate to
train on is rarely subjected to scrutiny: most studies train and
evaluate on the same surrogate without an out-of-loop fidelity check,
or evaluate on a different simulator without a clear bridge from
training-surrogate accuracy to deployment performance
\citep{Zhang2019WholeBuilding,Chen2018OptimalContext}. The
\emph{controller-side gap} the present paper addresses is the second:
existing DRL studies rarely separate the question of surrogate fidelity
from the question of controller performance.

Hierarchical decomposition has been proposed as one structural response
to long-horizon HVAC control problems, splitting setpoint planning
from low-level tracking
\citep{Bacon2017Options,Nachum2018HIRO,Liao2025HDRL}. Multi-objective
formulations \citep{Roijers2013MOSurvey,Yang2019MORL,Hayes2022MOSurvey}
are similarly motivated by the natural energy-comfort trade-off in
HVAC and have been used to expose tunable Pareto fronts at deployment
time \citep{Gao2024PredictiveFeatures,Sun2024StateHistory}. We adopt
both families as evaluation arms in Section~\ref{sec:results_control}.

\subsection{Surrogate and digital twin models for building energy}
\label{ssec:rw_surrogate}

Surrogate modelling for building energy is a mature subfield. Reviews
of data-driven approaches list dozens of architectures with reported
one-step and short-horizon RMSE numbers on benchmark testcases
\citep{Amasyali2018Review}. Recent contributions have moved from pure
black-box regression toward physically informed surrogates that retain
explicit thermodynamic parameters: low-order resistance--capacitance
(RC) models \citep{Bacher2011RCIdentification,Berthou2014RC,Bunning2020RC}
or RC--neural-ODE hybrids \citep{Chen2018NeuralODE,Chen2022PhysicalTwin}
calibrated by inverse procedures
\citep{Coraci2025StageCalibration,Wang2025SafetyMetric}. The recent
protocol of \cite{HouEvins2024} consolidates these efforts into a
four-stage development pipeline with explicit Reporting-Level-3
disclosure requirements; we adopt it as the methodological standard
for Block~1 of the present work.

A common feature of surrogate-fidelity studies, however, is that they
evaluate accuracy along trajectories that always start from a real
state. That setting is informative about \emph{predictive} accuracy
but only indirectly about the surrogate's value as a closed-loop RL
training environment, where every rollout is bootstrapped from the
surrogate's own output. \cite{Picard2017WhiteBox} highlights an
analogous distinction in the MPC context, and prior reviews of
surrogate-based MPC and RL emphasise the same point in passing
\citep{Mason2019RLReview,Drgona2020StochasticMPC}, but the failure mode
is rarely quantified directly. The \emph{fidelity-side gap} addressed
here is precisely that: showing the well-defined regime in which a
better predictive surrogate is not a better RL training environment.

\subsection{Physics-informed and physics-guided machine learning}
\label{ssec:rw_pinn}

The physics-informed machine-learning literature spans a wide range of
constructions, from hard-constraint physics-informed neural networks
(PINNs) \citep{Raissi2019PINN} through theory-guided data science
\citep{Karpatne2017TheoryGuided,Willard2022PIML} to neural ordinary
differential equations \citep{Chen2018NeuralODE}. Our construction is
deliberately on the soft-regularisation side: the physical model is
not enforced as a hard constraint, and it is not even part of the
policy's forward pass. It enters only as a penalty term that the
policy can violate at a cost. This is closer in spirit to auxiliary
tasks \citep{Jaderberg2017UNREAL} and Dyna-style world models
\citep{Sutton1991Dyna,Hafner2025Dreamer} than to PINNs, and we
position it as a controller-side regularisation strategy rather than
as a surrogate-architecture contribution.

\subsection{\boptest{} benchmarking}
\label{ssec:rw_boptest}

We adopt the \boptest{} benchmark \citep{Blum2021BOPTEST,
Arroyo2022BOPTESTService,BOPTESTrepo} as the reproducible testbed
throughout this study. Its \texttt{bestest\_air} testcase exposes a
single-zone air-side HVAC system with a documented comfort band, a
TMY weather file, and a composite safety/energy KPI ($m_s$) that
weights duration and severity of comfort violations against energy
use. Built-in PI control is available as a default reference; we use
it as the normaliser for Block~2 results
(Section~\ref{ssec:res2_pi}). Block~3 extends the evaluation to three
additional \boptest{} hydronic testcases under a pre-registered
manifest.

\subsection{Position of this work}
\label{ssec:rw_position}

This work sits at the intersection of the three gaps above. It
isolates the surrogate-fidelity question from the controller-performance
question through an explicit two-stage protocol (Block~1 then Block~2);
it shows the well-defined failure mode of fidelity-driven surrogates
as RL training environments
(Section~\ref{ssec:res1_failure_mode}); and it introduces hybrid
soft regularisation as the resolution, evaluated across three
controller families. Block~3 adds a pre-registered transferability
characterisation across three hydronic \boptest{} testcases and
reports the recipe's decomposable behaviour: the surrogate component
transfers structurally and the controller component does not, under
the chosen frozen-controller scope
(Section~\ref{sec:results_transfer}).
